%%%%%%%%%%%%%%%%%%%%%%%%%%%%%%%%%%%%%%%%%%%%%%%%%%%%%%%%%%%%%%%%%%%%%%%%%%%
% Handout und Uebungsblatt zu Lehreinheit 7
% Im Skript am Kapitelende, in handouts.tex als Druckvorlage.
%%%%%%%%%%%%%%%%%%%%%%%%%%%%%%%%%%%%%%%%%%%%%%%%%%%%%%%%%%%%%%%%%%%%%%%%%%%
\begin{handout}{Kognitive Verzerrungen}

\textbf{Worum es heute ging.} Drei gleich aussehende Schachteln, die
oberste mit ca.\ 50\,g, die anderen mit je ca.\ 3\,g. Allein gehoben
fühlt sich die oberste schwerer an als der ganze Stapel. Das Gehirn
erwartet ein Drittel des Gesamtgewichts und bekommt fast alles; die
enttäuschte Erwartung erzeugt die Illusion \parencite{charpentier1891}.
Sie bleibt, auch wenn Sie es wissen \parencite{flanagan2000independence}.

\textbf{Kernbegriffe}
\begin{description}
  \item[Bias] systematische, vorhersagbare Abweichung eines Urteils in
    eine Richtung \parencite{tversky1974judgment}.
  \item[Noise] unsystematische Streuung von Urteilen, die gleich ausfallen
    sollten \parencite{kahneman2021}.
  \item[Bestätigungsfehler] Belege so suchen, bewerten und erinnern, dass
    sie die eigene Meinung stützen.
  \item[Versuchsleitereffekt] Erwartungen der Forschenden verschieben die
    Daten \parencite{rosenthal1963experimenter}.
  \item[Hot-Cold-Empathy-Gap] im ruhigen Zustand unterschätzt man den
    Einfluss von Stress, Hunger, Erregung auf das eigene Urteil
    \parencite{loewenstein1996visceral}.
  \item[Rückschaufehler] wer das Ergebnis kennt, hält es für vorhersehbar
    \parencite{fischhoff1975hindsight}.
  \item[Publication Bias] Nullbefunde werden seltener veröffentlicht; die
    Literatur überschätzt Effekte \parencite{rosenthal1979filedrawer}.
\end{description}

\textbf{Verzerrung und Gegenmittel im Forschungsprozess}

\begin{center}
\small
\begin{tabular}{p{0.16\textwidth}p{0.33\textwidth}p{0.39\textwidth}}
\toprule
\textbf{Phase} & \textbf{typische Verzerrung} & \textbf{Gegenmittel} \\
\midrule
Planung & Bestätigungsfehler & Alternativhypothese, adversarielle Kollaboration \\
Erhebung & Versuchsleitereffekt, Anker & Blind- und Doppelblind-Design \\
Auswertung & p-Hacking, Forking Paths & Präregistrierung \\
Interpretation & Rückschaufehler, HARKing & Präregistrierung, Trennung geplant/explorativ \\
Publikation & Publication Bias & Registered Reports, Replikation \\
\bottomrule
\end{tabular}
\end{center}

\textbf{Replikationskrise:} Von 100 psychologischen Studien hatten
97\,\% signifikante Effekte berichtet, in der Replikation erreichten
etwa 36\,\% Signifikanz \parencite{opensciencecollaboration2015}.

\begin{merkbox}
Verzerrung ist der Normalbetrieb eines vorhersagenden Gehirns. Sie
verschwindet nicht, wenn man von ihr weiß. Bändigen lässt sie sich nur
durch Verfahren: blind messen, vorab festlegen, unabhängig urteilen.
\end{merkbox}

\textbf{Zum Weiterlesen:} \textcite{kahneman2011};
\textcite{nickerson1998ubiquitous}.
\end{handout}

\begin{uebungsblatt}{Kognitive Verzerrungen}

\aufgabe{Aufwärmen: ein fiktiver Kurzbericht}[in der Stunde, Einzelarbeit, 5 Min.]
\textit{Der folgende Bericht ist fiktiv und für diese Übung konstruiert.}

\begin{center}
\fbox{\parbox{0.92\textwidth}{\small
\textbf{\enquote{Barockmusik steigert die Konzentration}}\\[1mm]
(1)~Eine Dozentin, die seit Jahren Konzentrationstrainings mit Musik
anbietet, will zeigen, dass Barockmusik die Konzentration verbessert.
(2)~Sie gewinnt 24 Studierende aus ihrem eigenen Seminar. (3)~Alle lösen
zuerst am Vormittag Aufgaben ohne Musik und am Nachmittag desselben Tages
Aufgaben mit Musik. (4)~Die Dozentin bewertet die Lösungen selbst; auf
jedem Bogen ist vermerkt, ob er mit oder ohne Musik entstand. (5)~Erhoben
werden zwölf Konzentrationsmaße. Bei einem davon zeigt sich ein
signifikanter Vorteil der Musikbedingung; über die übrigen elf wird
nicht berichtet. (6)~Weil das Ergebnis nach 24 Personen \enquote{knapp
nicht signifikant} war, wurden zehn weitere Personen erhoben. (7)~In der
Diskussion heißt es, die Hypothese sei von Anfang an klar gewesen und
bestätige die langjährige Praxiserfahrung der Autorin. (8)~Eine frühere
Studie der Autorin ohne Effekt wurde nicht eingereicht.
}}
\end{center}

Notieren Sie zu jeder Nummer, ob dort ein Bias wirken könnte, welcher und
in welcher Phase.
\schreiblinien{5}

\aufgabe{Bias-Jagd im Artikel}[in der Stunde, Zweiergruppe, 25 Min.; Abschluss als Hausaufgabe]
Artikel (von der Lehrperson ausgewählt): \hrulefill

Finden Sie mindestens zwei Stellen, an denen ein Bias das Ergebnis
beeinflusst haben könnte. Prüfen Sie dafür nacheinander: Wer hat die
Studie durchgeführt und mit welchem Interesse? Wie wurden die
Teilnehmenden gewonnen und zugeteilt? Wer hat gemessen, und wusste diese
Person die Bedingung? Wie viele Messgrößen und Auswertungsvarianten gab
es, und wird über alle berichtet? Gab es eine Präregistrierung? Geht die
Interpretation über die Daten hinaus?

\begin{center}
\small
\renewcommand{\arraystretch}{2.6}
\begin{tabular}{|p{0.13\textwidth}|p{0.11\textwidth}|p{0.15\textwidth}|p{0.27\textwidth}|p{0.18\textwidth}|}
\hline
\textbf{Seite / Zitat} & \textbf{Phase} & \textbf{Bias} & \textbf{Begründung: in welche Richtung verschiebt er das Ergebnis?} & \textbf{Gegenmittel (verwendet?)} \\
\hline
 & & & & \\ \hline
 & & & & \\ \hline
 & & & & \\ \hline
\end{tabular}
\end{center}

Was haben die Autorinnen und Autoren methodisch gut gemacht?
\schreiblinien{2}

\aufgabe{Reflexion: der Bias in meinem eigenen Forschungsprozess}[Hausaufgabe]
Legen Sie Ihre Design-Skizze aus LE~6 zugrunde.

\begin{center}
\renewcommand{\arraystretch}{1.8}
\begin{tabular}{|p{0.36\textwidth}|p{0.54\textwidth}|}
\hline
Welcher Bias droht am stärksten? & \\ \hline
In welcher Phase? & \\ \hline
Woran würde ich merken, dass er gewirkt hat? & \\ \hline
Mit welchem \emph{Verfahren} begrenze ich ihn? & \\ \hline
\end{tabular}
\end{center}

\aufgabe{Tagebuch des Staunens}[Hausaufgabe]
Beschreiben Sie eine Situation, in der Sie etwas wussten und trotzdem
anders wahrgenommen oder entschieden haben. Was hätte Ihnen geholfen,
dem Wissen statt dem Erleben zu folgen?

\begin{lehrhinweis}[Erwartungshorizont]
\textbf{Aufgabe~1.} (1)~Interessenkonflikt, Bestätigungsfehler in der
Planung. (2)~Selektion und Abhängigkeit: eigene Studierende, mögliche
Gefälligkeitseffekte. (3)~Konfundierung der Bedingung mit Reihenfolge und
Tageszeit (Übung, Ermüdung); kein Bias im engeren Sinn, aber eine
Bedrohung der internen Validität, die in Richtung des erhofften Effekts
wirken kann. (4)~Versuchsleitereffekt, keine Blindbewertung (Erhebung).
(5)~selektives Berichten, p-Hacking über zwölf Maße (Auswertung).
(6)~optionales Weitererheben bis zur Signifikanz (Auswertung).
(7)~Rückschaufehler bzw.\ HARKing, keine Präregistrierung
(Interpretation). (8)~Publication Bias (Publikation).
\textbf{Aufgabe~2.} Mindestens zwei Stellen mit Fundort, benannter
Verzerrung, begründeter Richtung und konkretem Gegenmittel. Bloße
Schwächen (kleine Stichprobe, Gelegenheitsstichprobe) zählen als
Unsicherheit oder Einschränkung der externen Validität; als Bias nur, wenn
eine Richtung begründet wird.
\textbf{Aufgabe~3.} Ein Verfahren ist gefordert (Blindbewertung,
Präregistrierung, zweite Kodiererin, vorab festgelegtes
Abbruchkriterium), kein Vorsatz.
\end{lehrhinweis}
\end{uebungsblatt}
